In the previous two sections we proved the Main Theorem.  We showed that every compact surface in the boundary of a handlebody contains a simple closed curve that minimizes the girth over all closed curves in the surface.  In this section we describe how to construct such a curve.  This gives us an algorithm to compute the girth of any surface.  To describe the algorithm, we collect together some of the key constructions found in earlier sections of this paper.

Let $H$ be a handlebody, and let $S$ be a surface contained in 
$\bd H$.  Let $\dd$ be a disc system for $H$.  Recall from 
Section~\ref{introduction} that a disc system for $H$ is a disjoint
collection of discs that divides $H$ into a ball.  In 
Section~\ref{maintheorem}, we described a way to simplify
a disc system with respect to $S$.  Lemma~\ref{shortcut} describes 
the procedure.  It involves boundary-compressing along certain half-discs
until a locally minimal disc system is reached.
Let \locallyminimize{$\dd$} denote a disc system obtained 
from $\dd$ using this procedure.

Let $\dint{a}$ be a $\dd$-disc or a $\dd$-half-disc.  
Recall the associated diagram $T$.  Each region in $T$ is 
either segregated or desegregated.  
Let \segregation{$\dint{a}$} denote the number of regions in $T$ that
are segregated with respect to $\dint{a}$.

Let $n$ be \segregation{$\dint{a}$}.
When $\dint{a}$ is a $\dd$-disc (resp. $\dd$-half-disc),
Lemma~\ref{main lemma easy} (resp. Lemma~\ref{main lemma corollary})
describes a procedure to find a disc intersecting $S$ in no more
than $\ec{S}+n$ arcs (resp. $\ec{S}+\max(1,2n)$ arcs).
Let \constructdisc{$\dint{a}$} denote a disc obtained from $\dint{a}$ using 
this procedure.

Let $D$ be a disc intersecting $S$ in no more than $\ec{S}+n$ arcs.
Lemma~\ref{arcs in S} describes a procedure to find a simple
closed curve in $S$ that intersects $D$ in no more than $\max(1,n-1)$ 
points, or $n$ points if $\chi(S)$ is zero.  
Let \constructcurve{$D$} denote a simple closed curve obtained from $D$
using this procedure.

\subsection{The Algorithm.}
\label{surface algorithm}
Let $H$ be a handlebody, and let $S$ be a surface in $\partial H$.  
The algorithm below can be used to find a girth-minimizing curve in $S$
and a girth-realizing disc.

The algorithm takes as input a natural number $n$, a subsurface $S$,
and a disc system $\dd$.  It either terminates with no output or 
terminates after returning a simple closed curve $\gamma$ in $S$ and 
an essential disc $D$.  If it terminates with no output, then $\girth(S)>n$.
If it terminates with output, then $\girth(S)\leq n$.  The curve and disc 
satisfy $|\gamma\cap D|\leq\max(1,n)$.
Thus $\girth(\gamma)\leq n$.  The disc $D$ is girth-realizing for 
$\gamma$ unless they intersect once.
In this case a new disc can easily be formed that is disjoint from
$\gamma$ (Proposition~\ref{geometric girth one}).

To compute the girth, choose an arbitrary curve $\gamma$ in $S$.  
Compute the girth of $\gamma$ using the algorithm of 
Section~\ref{curvealgorithm}, say.  This establishes an upper bound
$n$ on the girth of $S$.  Repeat the algorithm below for successively
smaller $n$ until it terminates with no output.

The algorithm below is recursive in $\ec{S}$.  
The function \texttt{STOP} terminates all computation.


\bigskip

\begin{ttfamily}
Input: natural number $n$, surface $S$, disc system $\dd$

\smallskip

begin \GIRTH{$n$,$S$,$\dd$};

\hspace{20 pt}if ( $\chi(S)=2$ ) end GIRTH($n$,$S$,$\dd$);

\hspace{20 pt}$\dd\leftarrow$\locallyminimize{$\dd$};

\hspace{20 pt}for each $\dd$-disc $\dint{a}$;

\hspace{40 pt}if \segregation{$\dint{a}$}$\leq n$;

\hspace{60 pt}$D=$\constructdisc{$\dint{a}$};

\hspace{60 pt}$\gamma=$\constructcurve{$D$};

\hspace{60 pt}print $(\gamma, D)$;

\hspace{60 pt}STOP;

\hspace{40 pt}for each line of segregation $\ell$;

\hspace{60 pt}\GIRTH{$n$,$S-\ell$,$\dd$};

\hspace{20 pt}for each $\dd$-half-disc $\kappa$;

\hspace{40 pt}if \segregation{$\kappa$}$\leq\lfloor\frac{n}{2}\rfloor$;

\hspace{60 pt}$D=$\constructdisc{$\kappa$};

\hspace{60 pt}$\gamma=$\constructcurve{$D$};

\hspace{60 pt}print $(\gamma, D)$;

\hspace{60 pt}STOP;

\hspace{40 pt}for each line of segregation $\ell$;

\hspace{60 pt}\GIRTH{$n$,$S-\ell$,$\dd$};

\end{ttfamily}

\subsection{Proof of Algorithm.}

\begin{claim}
The algorithm of Section~\ref{surface algorithm} works and is effective.
\end{claim}

\begin{proof}
The algorithm terminates after finite time because at each stage there 
are finitely many $\dd$-discs, finitely many $\dd$-half-discs, and 
finitely many lines of segregation.  There are finitely many stages in the
recursion because cutting along $\ell$ reduces $\ec{S}$ by one.

To show the algorithm works, it suffices to show that if $\girth(S)\leq n$,
the algorithm terminates with output.  This is vacuously true for $\chi(S)=2$.
So assume we have shown it for $\chi(S)\geq k+1$.  We will prove
it for $\chi(S)=k$.

Assume that $\girth(S)\leq n$.  
Let $\dd$ be a locally minimal disc system.
Let $\gamma$ be a closed curve in $S$ with girth $n$.  Let $D$ be a 
girth-realizing disc for $\gamma$.  Apply the Subordination Lemma to $D$.

There is either a $\dd$-disc subordinate to $D$, or two $\dd$-half-discs 
independently subordinate to $D$.  Depending on the case, let $\dint{a}$ 
be that $\dd$-disc, or the $\dd$-half-disc with smaller complexity.
Let $m$ be the complexity of $\dint{a}$.  Then $\gamma$ homotopically
intersects $\dint{a}$ exactly $n$ times.  If $\dint{a}$ is the $\dd$-disc
(resp. $\dd$-half-disc), then $m\leq n$ (resp. 
$m\leq\lfloor\frac{n}{2}\rfloor$).

By Lemma~\ref{segregated regions}, either \segregation{$\dint{a}$}$\leq m$, 
or $\gamma$ is isotopic to a curve disjoint from some line of segregation.  
Note that if $\chi (S)=0$, the first must be true.  In the former case, 
an application of \texttt{CONSTRUCT.DISC} and \texttt{CONSTRUCT.CURVE} 
finds a curve $\gamma$ and disc $D$ satisfying $|\gamma\cap\dd|\leq\max(1,n)$.
In the latter case, a girth-minimizing curve is contained 
in $S-\ell$ for some line of segregation 
(see the proof of Lemma~\ref{segregated regions}).
Then apply \texttt{GIRTH} to $S-\ell$, whose Euler characteristic 
is $\chi(S)=k+1$.
\end{proof}













